\documentclass{../../templates/karthicv-single}
\usepackage{color}
\setlength{\parindent}{0pt}

% RAEng Global Talent CV - peer review route
\renewcommand{\cvname}{Karthi Arunachalam}
\renewcommand{\cvemail}{k18arma@gmail.com}
\renewcommand{\cvphone}{+44 7859 745670}
\renewcommand{\cvheadline}{Engineering researcher and emerging climate technology founder in chemical and process engineering}
\renewcommand{\cvwebpage}{https://www.linkedin.com/in/k8r/}
\renewcommand{\cvaddress}{Cambridge, UK}

\begin{document}

\maketitle

\vspace{-1em}
\begin{center}
\textit{Chemical and process engineering researcher working across carbon capture, advanced materials, energy storage, automation, and computational engineering. Experienced in building experimental systems, reducing uncertainty around emerging climate technologies, and translating technical insight into patents, industrial validation, confidential customer-facing research outputs, and venture formation.}
\end{center}

\section{Research and Engineering Experience}

\experience{Founder}{ExerVolt}{Cambridge, UK}{9/25-Present}
{Building an emerging energy storage and conversion venture focused on long-duration energy storage.\\
Developing engineering and computational capabilities that connect materials, electrochemical performance, manufacturing, techno-economics, and deployment strategy.\\
Built early digital twin and techno-economic models to evaluate beachhead markets in the UK and India.\\
Working with Carbon13 in Cambridge to refine product-market fit, technical roadmap, and commercialisation strategy.}

\experience{Senior Research Scientist}{SLB, New Energy Research}{Cambridge, UK}{9/21-Present}
{Led climate and energy technology research projects with full project lifecycle responsibility, including technical planning, experimental delivery, stakeholder management, and reporting.\\
Established a carbon capture characterisation laboratory at SLB Cambridge in under six months, enabling repeatable validation and comparative assessment of next-generation sorbents and carbon capture configurations.\\
Led technical de-risking of Immaterial's monolithic metal organic framework carbon capture platform, combining experimental benchmarking, modelling insight, and structured performance validation under industrially relevant conditions. This work materially informed SLB's confidence in the technology and contributed to its participation in Immaterial's Series A2 funding round.\\
Worked with TWI on a high pressure CO2 permeation monitoring system for subsea carbon sequestration infrastructure, integrating electrochemical sensing, hardware redesign, automation, signal processing, and modelling-led optimisation for operation around 200 bar and 60 degC.\\
Lead author of confidential SLB research notes on CO2 measurements for Petrobras/TWI riser materials testing and NiH2 battery modelling.\\
Developed computational and data-driven engineering workflows, including graph neural network based techno-economic assessment, Bayesian uncertainty analysis, digital twins, and automated laboratory data pipelines.\\
Managed cross-functional technical work across materials, process, mechanical, electrical, automation, and commercial stakeholders.}

\experience{Chemical Engineer}{Exactmer}{London, UK}{2/21-9/21}
{Developed automation and control systems for pharmaceutical oligonucleotide synthesizers.\\
Designed and built membrane testing automation using programmable control, automated fluid handling, and real-time data acquisition.\\
Designed process improvements for solvent recovery and manufacturing carbon footprint reduction.\\
Participated in HAZOP studies as automation and process safety contributor.}

\experience{PhD Researcher}{University of Edinburgh}{Edinburgh, UK}{9/16-2/21}
{Conducted research on sorbent materials and reaction-diffusion systems for CO2 sequestration and carbon capture.\\
Designed, built, and commissioned high pressure, high temperature multi-component gas sorption systems.\\
Developed experimental and numerical workflows for gas sorption, transport phenomena, and carbon storage analysis.\\
Published peer-reviewed and conference research, and taught undergraduate Mass Transfer.}

\section{Commercial Outputs and Engineering Impact}

\textbf{Filed patent applications}: Four filed United States patent applications in carbon capture, advanced materials, and chemical processing systems:\\
US Serial No. 19/395,010, ``Multi-Pellet Packed Bed Technology'', filed November 20, 2025.\\
US Serial No. 19/395,050, ``Synthesis of Large-Scale Monolithic Structures and Thin-Film Structures that Function as Sorbents or Catalysts'', filed November 20, 2025.\\
US Serial No. 19/395,067, ``Deposition of Adsorbent Material and/or Catalyst Material on Metal Surfaces of a Chemical Processing Apparatus'', filed November 20, 2025.\\
US Serial No. 19/361,338, ``Systems and Methods for Metal Organic Framework (MOF) Stack Kinetics'', filed October 17, 2025.\\
\textbf{Industrial de-risking}: Led engineering validation work that materially informed SLB's confidence in Immaterial's carbon capture technology and contributed to SLB's participation in Immaterial's Series A2 funding round.\\
\textbf{Laboratory creation}: Built a functional carbon capture characterisation laboratory enabling repeatable testing and comparison of next-generation sorbents.\\
\textbf{Venture formation}: Building ExerVolt, a pre-seed-stage venture developing UK-based technical capability and intellectual property in long-duration energy storage and energy conversion.\\

\section{Education and Professional Qualifications}

\cventry{Certificate of Achievement in Machine Learning with Graphs}{Stanford Online}\\
\textit{Graph neural networks, representation learning, knowledge graphs, node embeddings, graph generative models, and large-scale graph analysis} \hfill 2024\\
\cventry{Battery Modelling Certificate}{MathWorks / MATLAB Academy}\\
\textit{Computational battery modelling and simulation for electrochemical energy systems} \hfill 2024\\
\cventry{PhD Engineering}{University of Edinburgh}\\
\textit{Sorbent materials and reaction-diffusion systems for CO2 sequestration and carbon capture} \hfill 2021\\
\cventry{MSc Management}{University of Aberdeen}\\
\textit{Project management and strategy; dissertation on financial modelling and carbon pricing for CCS} \hfill 2015\\
\cventry{BEng Chemical Engineering}{Anna University}\\
\textit{First Class Honours; process control systems} \hfill 2013\\

\section{Selected Research Outputs}

K. Arunachalam and X. Fan, ``Impact of heat of sorption on thermal enhanced recovery of sorbed gas from gas shale reservoirs - an experimental and simulation study'', \textbf{Journal of Natural Gas Science and Engineering}, 2020.\\
K. Arunachalam and X. Fan, ``Implications of Sorption on Carbon Dioxide Sequestration in Gas Shales'', \textbf{COMSOL International Conference}, 2020.\\
K. Arunachalam (lead author), E. Corcoran, A. Meredith, and D. Campos de Faria, ``Carbon Dioxide Measurements for Petrobras/TWI Permeance Testing of Riser Materials'', \textbf{SLB Cambridge Research Note}, Chemistry, 2022. Prepared for Petrobras and TWI. Confidential, restricted distribution.\\
K. Arunachalam (lead author), A. Meredith, D. Campos de Faria, and J. Rubio, ``NiH2 Battery Modelling'', \textbf{SLB Cambridge Research Note}, Energy Storage, 2023. Confidential, restricted distribution.

\section{Professional Membership and Service}

\experience{Associate Member}{Institution of Chemical Engineers (IChemE)}{UK}{Current}
{Professional body membership in chemical engineering.}

\experience{Committee Member}{IChemE Clean Energy Special Interest Group}{UK}{7/23-Present}
{Contributed to clean energy community activity, technical discussion, and professional engagement around sustainable engineering applications.}

\experience{Secretary}{IChemE Scottish Members Group}{Edinburgh, UK}{5/17-6/19}
{Coordinated committee activity, meetings, AGMs, and member engagement for the regional chemical engineering community.}

\experience{Tutor}{University of Edinburgh}{Edinburgh, UK}{9/16-9/20}
{Taught undergraduate Mass Transfer in classroom and laboratory settings.}

\section{Selected Technical Skills}

\textbf{Chemical and process engineering}: carbon capture, gas sorption, transport phenomena, CO2 monitoring, process design, HAZOP, process safety, experimental systems.\\
\textbf{Energy and climate technology}: long-duration energy storage, battery modelling, techno-economic assessment, industrial decarbonisation, carbon capture and storage.\\
\textbf{Computational engineering}: Python, MATLAB, COMSOL, Bayesian modelling, graph neural networks, digital twins, automation, data acquisition, signal processing.\\
\textbf{Technical leadership}: cross-functional project delivery, external collaboration, stakeholder management, laboratory build-out, technology de-risking, venture formation.

\end{document}
